\section{Wielandt Bound}
\label{ch:wielandt_supporting_results}
This section collects exact-word, cumulative-span, augmentation, blocking,
complementary-gap, and one-step padding results used by the main
quantitative argument earlier in this chapter.

\section{Exact-word and block-injectivity support}
\label{sec:app_wielandt_exact_word_support}

\begin{lemma}[Two-sided span from one nonzero matrix]
    \label{lem:two_sided_nonzero_matrix_span}
    \lean{Matrix.span_range_mul_nonzero_mul_eq_top}
    \leanok
    If $C\in\MN{D}$ is nonzero, then
    \begin{align}
        \spn_{\C}\{RCS:R,S\in\MN{D}\}
         & =\MN{D}.
        \label{eq:wld_two_sided_nonzero_span}
    \end{align}
    This is the matrix-factorization input used in~%
    \cite[Lemma~3]{PerezGarcia2007Matrix}.
\end{lemma}

\begin{proof}\leanok
    Choose a nonzero entry $C_{pq}$. Then every matrix unit $E_{ij}$
    is a scalar multiple of $E_{ip}CE_{qj}$, so the span in
    (\ref{eq:wld_two_sided_nonzero_span}) contains the standard matrix basis.
\end{proof}

\section{Cumulative-span and spectral linear algebra}
\label{sec:app_wielandt_cumulative_support}

The next two lemmas connect irreducible matrix actions with the cumulative word
span used throughout Chapter~\ref{ch:wielandt}.

\begin{definition}[Fitting decomposition]\label{def:fitting_decomp}
    \lean{Wielandt.FittingDecomposition}
    \leanok
    A \emph{Fitting decomposition} of a linear endomorphism
    $f:V\to V$ on a finite-dimensional vector space over an
    algebraically closed field consists of:
    \begin{enumerate}
        \item $f$ is nilpotent on the generalized $0$-eigenspace,
        \item $f$ is invertible on each generalized $\mu$-eigenspace
              for $\mu\neq0$,
        \item the generalized eigenspaces span~$V$,
        \item the generalized eigenspaces are linearly independent.
    \end{enumerate}
\end{definition}

\begin{theorem}[Fitting decomposition exists]\label{thm:fitting_decomp}
    \lean{Wielandt.fittingDecomposition}
    \leanok
    \uses{def:fitting_decomp}
    Every linear endomorphism on a finite-dimensional vector space
    over an algebraically closed field admits a Fitting decomposition.
\end{theorem}

\begin{proof}
    \leanok
    Nilpotency on the zero generalized eigenspace follows from
    the definition of generalized eigenspaces.
    Invertibility on nonzero generalized eigenspaces follows
    because $f-\mu$ is nilpotent there, so
    $f=\mu(1-(1-f/\mu))$ is invertible.
    Spanning and independence are standard results for
    generalized eigenspaces over algebraically closed fields.
    In~\cite{Sanz2010Quantum} and \cite[Lemma~6.3(b)]{Wolf2012Quantum},
    the Jordan normal form is used directly.
    The argument is phrased in terms of generalized eigenspaces instead.
\end{proof}

\begin{theorem}[Nilpotent power bound]\label{thm:nilpotent_pow_bound}
    \lean{Wielandt.nilpotent_pow_eq_zero_of_finrank}
    \leanok
    If $f$ is nilpotent on a space of dimension~$n$, then $f^n=0$.
\end{theorem}

\begin{proof}\leanok
    A nilpotent endomorphism on an $n$-dimensional space has
    nilpotency index at most~$n$.
\end{proof}

\begin{theorem}[Nonzero trace implies eigenvector]\label{thm:eigenvector_from_trace}
    \lean{exists_eigenvector_of_trace_ne_zero}
    \leanok
    If $M\in\MN{D}$ has $\tr(M)\neq0$, then $M$ has a nonzero eigenvalue
    $\mu\neq0$ and a corresponding eigenvector $\varphi\neq0$ satisfying
    $M\varphi=\mu\varphi$.
\end{theorem}

\begin{proof}
    \leanok
    The trace is the sum of the eigenvalues of $M$, counted with
    multiplicity, so a nonzero trace gives a nonzero eigenvalue~$\mu$.
    Choosing a nonzero vector $\varphi\in\ker(M-\mu\Id)$ gives
    $M\varphi=\mu\varphi$.
\end{proof}

\section{One-step augmentation}
\label{sec:app_wielandt_augmentation}

\section{Blocking and fixed-length spanning}
\label{sec:app_wielandt_blocking_support}
\label{sec:rectangular_span}

\begin{theorem}[Blocking identities for exact word spans]
    \label{thm:kraus_blocking_word_span}
    \lean{Kraus.evalWord_blockTensor,
        Kraus.wordSpan_blockTensor_one,
        Kraus.wordSpan_blockTensor,
        Kraus.HasEventuallyFullWordSpan.blockTensor}
    \leanok
    \uses{def:kraus_eval_word}
    Fix a block length $L$ and regard the words of length $L$ in a Kraus
    family $K$ as a new family $K^{[L]}$. Evaluation of a blocked word agrees
    with evaluation of its flattened word, and consequently
    \begin{align}
        K^{[L]}_n=K_{nL}.
        \notag
    \end{align}
    If $L>0$ and the exact word spans of $K$ are eventually full, then the
    exact word spans of $K^{[L]}$ are eventually full as well.
\end{theorem}

\begin{proof}\leanok
    The evaluation identity follows by induction on the blocked word. Taking
    spans gives the displayed equality. For eventual fullness, sufficiently
    large $n$ also makes $nL$ exceed the original threshold.
\end{proof}

\section{Complementary-gap consequences}
\label{sec:app_wielandt_complementary_gap}

\section{Identity in the one-step span}
\label{sec:app_wielandt_identity_in_span}

\begin{remark}[Counterexample: cumulative spanning does not imply normality]%
    \label{rem:cumspan_counterexample}\leanok
    In $\MN{2}$, let $A^0=E_{12}$ and $A^1=E_{21}$ be the off-diagonal
    matrix units.
    Then $\algspn(A)=\MN{2}$ and $T_2(A)=\MN{2}$, but the word spans
    alternate: $S_n(A)$ contains only off-diagonal matrices for odd~$n$
    and only diagonal matrices for even~$n$.
    Hence $S_n(A)\neq\MN{2}$ for every~$n$, and $A$ is not normal.
    The one-step padding condition fails because
    $\Id\notin S_1(A)=\spn_{\C}\{E_{12},E_{21}\}$.
\end{remark}
